%\RequirePackage{kvoptions-patch}
\documentclass[a4paper, twoside, bibliography=totoc, headsepline, cleardoublepage=empty, parskip=half, draft=false]{scrbook}

\let\ifdeutsch\iffalse
\let\ifenglish\iftrue

\input{config}


\usepackage[
  title={
Ein Spiel Zur Erkennung von Deepfakes Zur Sensibilisierung Der Nutzer Durch Öffentliche Displays}, % Do not forget to capitalize your title correctly, you may use the following page to help you: https://capitalizemytitle.com/
  author={Anton Hantschmann},
  orcid=0009-0000-6440-2891, % get your own ORCID via https://orcid.org/
  email={anton.hantschmann@campus.lmu.de},
  type=bachelor,
  institute={Institut für Informatik},
  course={Medieninformatik},
  examiner={Prof.\ Dr.\ Florian Alt},
  supervisor={Doruntina Murtezaj,\ M.Sc.},
  startdate={May 13, 2025},
  enddate={August 5, 2025},
  copyright=ccbysa
  language=german,
  aids=true, % Spefici if you used any aids in the crated of the thesis, if yes [true] if no adis were used [false]
  % Next, enter which tools were used and how according to the following taxonomy (https://ceur-ws.org/GenAI/Taxonomy.html). If no aid was used, the next three values can be ignored.
  aidstools={Claude}, 
  adiscontributions={Grammar and spelling check, Paraphrase and reword, Peer review simulation},
  aidsextracontributions={Claude Sonnet wurde zum debugging und Fehler im Code finden genutzt.}
]{lmu-thesis-cover}

\input{acronyms}

\makenoidxglossaries

\begin{document}

\input{commands}

\pagenumbering{arabic}
\Coverpage
\Copyright
%Eigener Seitenstil fuer die Kurzfassung und das Inhaltsverzeichnis
\deftriplepagestyle{preamble}{}{}{}{}{}{\pagemark}
%Doku zu deftriplepagestyle: scrguide.pdf
\pagestyle{preamble}
\renewcommand*{\chapterpagestyle}{preamble}



%Kurzfassung / abstract
%auch im Stil vom Inhaltsverzeichnis
\section*{Kurzfassung}
Die zunehmende Verbreitung und Verbesserung von Deepfake-Technologien stellt ein wachsendes gesellschaftliches Risiko dar, 
während der Bevölkerung gleichzeitig die notwendigen Fähigkeiten zur Erkennung manipulierter Medien fehlen.
Diese Bachelorarbeit entwickelte daher ein interaktives Lernspiel für Public Displays zur Schulung von Deepfake-Erkennungsfähigkeiten. 
Die Anwendung kombinierte echte und manipulierte Videos mit Erkennungsstrategien sowie spielerischen Elementen wie Fortschrittsverfolgung und Performance-Statistiken.
Eine zweiphasige Feldstudie evaluierte die Wirksamkeit in öffentlicher Umgebung. 
Der spielbasierte Ansatz erwies sich als effektiv für die Vermittlung sicherheitskritischer Themen, 
wobei die Teilnehmer Verbesserungen ihrer Selbsteinschätzung zur Deepfake Erkennung und hohe Lernmotivation zeigten. 
Public Displays als Lernplattform fanden große Akzeptanz.
Besonders bemerkenswert war die soziale Komponente: 
Das öffentliche Setting regte spontane Diskussionen zwischen Nutzern über Deepfake-Problematiken an.
Die Studie zeigt, dass Public Display-basierte Spiele nicht nur Bewusstsein für Sicherheitsthemen schaffen, sondern auch gesellschaftliche Diskurse anstoßen.
Dies eröffnet neue Perspektiven für öffentliche Bildungsmaßnahmen im Bereich sicherheitskritischer Themen.

\cleardoublepage

\section*{Abstract}
The increasing spread and improvement of deepfake technologies poses a growing social risk, 
while at the same time the population lacks the necessary skills to recognize manipulated media.
This bachelor's thesis therefore developed an interactive educational game for public displays to train deepfake recognition skills. 
The application combined real and manipulated videos with recognition strategies and playful elements such as progress tracking and performance statistics.
A two-phase field study evaluated its effectiveness in public settings. 
The game-based approach proved effective for teaching security-critical topics, 
with participants showing improvements in their self-assessment of deepfake detection and high learning motivation. 
Public displays as a learning platform were well accepted.
Particularly noteworthy was the social component: 
the public setting stimulated spontaneous discussions among users about deepfake issues.
The study shows that public display-based games not only raise awareness of security issues, but also spark social discourse.
This opens up new perspectives for public education measures in the area of security-critical topics.

\cleardoublepage


% BEGIN: Verzeichnisse

\iftex4ht
\else
  \microtypesetup{protrusion=false}
\fi

%%%
% Literaturverzeichnis ins TOC mit aufnehmen, aber nur wenn nichts anderes mehr hilft!
% \addcontentsline{toc}{chapter}{Literaturverzeichnis}
%
% oder zB
%\addcontentsline{toc}{section}{Abkürzungsverzeichnis}
%
%%%

%Produce table of contents
%
%In case you have trouble with headings reaching into the page numbers, enable the following three lines.
%Hint by http://golatex.de/inhaltsverzeichnis-schreibt-ueber-rand-t3106.html
%
%\makeatletter
%\renewcommand{\@pnumwidth}{2em}
%\makeatother
%
\tableofcontents

% Bei einem ungünstigen Seitenumbruch im Inhaltsverzeichnis, kann dieser mit
% \addtocontents{toc}{\protect\newpage}
% an der passenden Stelle im Fließtext erzwungen werden.

\listoffigures
\listoftables

% Control List of Listings
\let\iflistings\iffalse
%Wird nur bei Verwendung von der lstlisting-Umgebung mit dem "caption"-Parameter benoetigt
%\lstlistoflistings
%ansonsten:
\iflistings
  \ifdeutsch
    \listof{Listing}{Verzeichnis der Listings}
  \else
    \listof{Listing}{List of Listings}
  \fi
\fi

% Control List of Algorithms
\let\ifalgorithms\iffalse
\ifalgorithms
  %mittels \newfloat wurde die Algorithmus-Gleitumgebung definiert.
  %Mit folgendem Befehl werden alle floats dieses Typs ausgegeben
  \ifdeutsch
    \listof{Algorithmus}{Verzeichnis der Algorithmen}
  \else
    \listof{Algorithmus}{List of Algorithms}
  \fi
  %\listofalgorithms %Ist nur für Algorithmen, die mittels \begin{algorithm} umschlossen werden, nötig
\fi

% Control Glossary
\let\ifglossary\iffalse
\ifglossary
    \ifdeutsch
        \printnoidxglossary[type=\acronymtype, title=Abkürzungsverzeichnis]
    \else
        \printnoidxglossary[type=\acronymtype, title=List of Acronyms]
    \fi
\fi

\iftex4ht
\else
  %Optischen Randausgleich und Grauwertkorrektur wieder aktivieren
  \microtypesetup{protrusion=true}
\fi
% END: Verzeichnisse


% Headline and footline
\renewcommand*{\chapterpagestyle}{scrplain}
\pagestyle{scrheadings}
\pagestyle{scrheadings}
\ihead[]{}
\chead[]{}
\ohead[]{\headmark}
\cfoot[]{}
\ofoot[\usekomafont{pagenumber}\thepage]{\usekomafont{pagenumber}\thepage}
\ifoot[]{}


%% vv  scroll down for content  vv %%
\input{content.tex}


%TC:ignore
\printbibliography

All links were last followed on \today{}.

\newpage
\appendix
\input{survey_questions.tex}
%\input{latexhints/latexhints-english}

\Aids

\Affirmation
%TC:endignore
\end{document}